% MASTER 5 -- CMOS pseudo-static D flip-flop using transmission gates (Q4).
% Master TG (phi) samples D into an inverter; slave TG (phi-bar) closes the
% feedback inverter loop to latch/hold (the feedback "refreshes" -> pseudo-static).
\documentclass[border=10pt]{standalone}
\usepackage{tikz}
\usetikzlibrary{arrows.meta}
\begin{document}
\begin{tikzpicture}[>=Latex, font=\small,
  tg/.style={draw, thick, trapezium, trapezium angle=70, minimum width=0.9cm}]
\usetikzlibrary{shapes.geometric}
% input inverter
\node[draw, thick, isosceles triangle, isosceles triangle apex angle=60,
      minimum width=1cm] (inv0) at (0,0) {};
\node at (-0.2,0) {\scriptsize $\triangleright$};
\draw[->] (-2.0,0) node[left]{$D$} -- (inv0.west);
% master TG (phi)
\node[draw, thick, minimum width=1.1cm, minimum height=0.8cm] (tg1) at (2.6,0) {TG};
\node[above, font=\scriptsize] at (2.6,0.45) {$\phi$};
\draw[->] (inv0.east) -- (tg1.west);
% storage / output inverters (buffer + feedback)
\node[draw, thick, isosceles triangle, isosceles triangle apex angle=60,
      minimum width=1cm] (inv1) at (4.8,0) {};
\node at (4.6,0) {\scriptsize $\triangleright$};
\draw[->] (tg1.east) -- (inv1.west) node[midway,above,font=\scriptsize]{};
\node[draw, thick, isosceles triangle, isosceles triangle apex angle=60,
      minimum width=1cm] (inv2) at (7.2,0) {};
\node at (7.0,0) {\scriptsize $\triangleright$};
\draw[->] (inv1.east) -- (inv2.west);
\draw[->] (inv2.east) -- (9.2,0) node[right]{$Q$};
% slave TG (phi-bar) closing the feedback loop
\node[draw, thick, minimum width=1.1cm, minimum height=0.8cm] (tg2) at (6.0,-1.8) {TG};
\node[below, font=\scriptsize] at (6.0,-2.25) {$\overline{\phi}$};
\draw (inv2.east) -- (8.4,0) -- (8.4,-1.8) -- (tg2.east);
\draw[->] (tg2.west) -- (4.0,-1.8) -- (4.0,-0.0) -- (inv1.west);
\end{tikzpicture}
\end{document}
